%\yxy{add more details to the figure to make it more interesting, like example systems (only the most important ones)}}
% new feedback, make it more colorful, instead of showing example systems, show a short description
\begin{figure}
    \centering
    \begin{tikzpicture}[
        font=\small\sffamily,
        cell/.style={draw, rounded corners=2pt, minimum width=2.5cm,
        minimum height=1.4cm, align=center,
        font=\small\sffamily\bfseries},
        open/.style={cell, dashed, line width=0.9pt, fill=black!8},
        ax/.style={font=\small\sffamily\bfseries, align=center},
        flow/.style={-{Latex[length=2.2mm]}, thick, black!65},
    ]

        % columns: relational, vector; rows: OLTP, CPU OLAP, GPU OLAP
        \node[cell] (r1) at (0,0)      {70s--80s\\[2pt] {\scriptsize\mdseries Oracle, PostgreSQL,}\\ {\scriptsize\mdseries MySQL, SQL Server}};
        \node[cell] (r2) at (0,-2.5)   {$\sim$mid-1990s\\[2pt] {\scriptsize\mdseries Snowflake, BigQuery,}\\ {\scriptsize\mdseries Databricks, DuckDB}};
        \node[cell] (r3) at (0,-5.0)   {2010s\\[2pt] {\scriptsize\mdseries Sirius, Theseus,}\\ {\scriptsize\mdseries HeavyDB, Kinetica}};
        \node[cell] (v1) at (4.5,0)    {late 2010s\\[2pt] {\scriptsize\mdseries Milvus, Pinecone,}\\ {\scriptsize\mdseries pgvector, Oracle}};
        \node[cell] (v2) at (4.5,-2.5) {2024\\[2pt] {\scriptsize\mdseries Databricks, Spark,}\\ {\scriptsize\mdseries BigQuery, DuckDB}};
        \node[open] (v3) at (4.5,-5.0) {this work\\[2pt] {\scriptsize\mdseries Sirius}};

        % arrows
        \draw[flow] (r1) -- node[fill=white, inner sep=1pt, font=\scriptsize\sffamily\mdseries, align=center] {data warehousing,\\column stores} (r2);
        \draw[flow] (r2) -- node[fill=white, inner sep=1pt, font=\scriptsize\sffamily\mdseries, align=center] {memory bandwidth,\\cost efficiency} (r3);
        \draw[flow] (v1) -- node[fill=white, inner sep=1pt, font=\scriptsize\sffamily\mdseries] {batch ML workloads} (v2);
        \draw[flow] (v2) -- node[fill=white, inner sep=1pt, font=\scriptsize\sffamily\mdseries] {compute throughput} (v3);
        \draw[flow] (r1) -- node[above, font=\scriptsize\sffamily\mdseries, align=center] {embedding\\models} (v1);
        \draw[flow] (r3) -- node[above, font=\scriptsize\sffamily\mdseries, align=center] {embeddings in\\the warehouse} (v3);

        % column axis (x)
        \node[ax, anchor=south] at ([yshift=3pt]r1.north) {Relational};
        \node[ax, anchor=south] at ([yshift=3pt]v1.north) {Vector};

        % row axis (y)
        \node[ax, anchor=east, inner xsep=0pt] at ([xshift=-6pt]r1.west) {OLTP};
        \node[ax, anchor=east, inner xsep=0pt] at ([xshift=-6pt]r2.west) {CPU\\OLAP};
        \node[ax, anchor=east, inner xsep=0pt] at ([xshift=-6pt]r3.west) {GPU\\OLAP};
    \end{tikzpicture}
    \caption{
        Each cell shows when the transition began and its representative systems (not necessarily from that period), and arrows label what drove it. Relational data moved from transactional to analytical processing on CPUs~\cite{StonebrakerPavloAround24}, then to GPUs~\cite{AramburuTheseus25, HeavyDB, Kinetica, Yogatama26}. Vector data is following the same path: vector search serves point queries~\cite{StonebrakerPavloAround24, WangMilvus21}, and the vector join is now a SQL clause in several CPU engines~\cite{DuckDBVSSBlog24, DatabricksNearestByJoinBlog26, DuckDBNearestByPR24137, SparkNearestByJoin26}. The vector join on the GPU is the open cell that this work fills.
    }
    \label{fig:quadrant}
\end{figure}
